\chapter{Quantum field theories as functors (KS Section 3)}

This chapter is the blueprint for Section~3 of Kontsevich--Segal,
\emph{Wick rotation and the positivity of energy in quantum field theory}
(arXiv:2105.10161), to be formalized in \texttt{KontsevichSegal/Cobordism/}
and \texttt{KontsevichSegal/FieldTheory/}. The section recasts a
$d$-dimensional field theory as a holomorphic functor on a category
$\mathcal C_d^{\mathbb C}$ of space-time manifolds carrying complex metrics.

\section{The cobordism category}

\begin{definition}[The cobordism category $\mathcal C_d^{\mathbb C}$; KS Section 3]
  \label{def:cobordism-category}
  \notready
  \uses{def:allowable}
  A \emph{complex metric} on a smooth manifold $X$ is a smooth field
  $x \mapsto g_x$ assigning to each point an allowable complex metric
  $g_x \in Q_{\mathbb C}(T_x X)$ on the tangent space. Write
  $\mathrm{Met}_{\mathbb C}(X)$ for the space of such fields.

  The category $\mathcal C_d^{\mathbb C}$ is defined as follows.

  \emph{Objects.} An object is a germ of a $d$-manifold along a closed
  two-sided $(d-1)$-manifold $\Sigma$. Concretely, $\Sigma$ is a compact
  $(d-1)$-manifold presented as a closed submanifold of a $d$-manifold $U$,
  and any two open neighbourhoods of $\Sigma$ in $U$ define the same object.
  The submanifold $\Sigma$ is two-sided in $U$ and carries a
  \emph{co-orientation}, a choice of which of its two sides is incoming and
  which is outgoing. The germ carries a complex metric, that is a germ along
  $\Sigma$ of an element $g \in \mathrm{Met}_{\mathbb C}(U)$. One usually
  suppresses the thickening $U$, the co-orientation, and the metric $g$ from
  the notation.

  \emph{Morphisms.} A morphism $\Sigma_0 \to \Sigma_1$ is a cobordism
  $M : \Sigma_0 \leadsto \Sigma_1$, a compact $d$-manifold carrying a complex
  metric $g_M \in \mathrm{Met}_{\mathbb C}(M)$, whose incoming and outgoing
  boundary germs are the objects $\Sigma_0$ and $\Sigma_1$ with their
  co-orientations and metric germs. Two cobordisms
  $M, M' : \Sigma_0 \leadsto \Sigma_1$ represent the same morphism when there
  is an isometry $M \to M'$ restricting to the identity on the germs
  $\Sigma_0$ and $\Sigma_1$.

  \emph{Composition.} Composition is concatenation of cobordisms. A cobordism
  $\Sigma_0 \leadsto \Sigma_1$ followed by a cobordism
  $\Sigma_1 \leadsto \Sigma_2$ is glued along their common boundary germ
  $\Sigma_1$ to give a cobordism $\Sigma_0 \leadsto \Sigma_2$. Taking objects
  to be germs rather than bare $(d-1)$-manifolds is exactly what makes the
  smooth structure on the concatenation well defined, since there is no
  canonical smooth structure on the concatenation of two smooth cobordisms.

  \emph{Identity morphisms.} $\mathcal C_d^{\mathbb C}$ is not literally a
  category, as it has no identity morphisms. The paper records two ways to
  deal with this, neither resolved here. One may declare an isomorphism
  $\Sigma_0 \to \Sigma_1$ to be a cobordism of zero length; such degenerate
  cobordisms exist, but in any field theory they are represented by operators
  that are not nuclear, so they fall outside the intended class of morphisms.
  The alternative, which the paper adopts, takes a field theory to be a
  functor $\Sigma \mapsto E_{\Sigma}$ from $(d-1)$-manifold germs and their
  isomorphisms to vector spaces, together with a transformation
  $Z_M : E_{\Sigma_0} \to E_{\Sigma_1}$ for each cobordism. The true
  replacement for the missing identity operators is then a condition imposed
  on field theories, namely that the canonical maps
  $\check E_{\Sigma} \to \hat E_{\Sigma}$ are injective with dense image.

  Stating this in Lean requires smooth-manifold infrastructure not developed
  here: germs of manifolds along a submanifold, two-sidedness and
  co-orientation, isometries of manifolds carrying complex metrics, and gluing
  of cobordisms. It also requires realizing $\mathrm{Met}_{\mathbb C}$ as the
  sections of a bundle whose fibre at $x$ is the domain $Q_{\mathbb C}(T_x X)$
  of allowable complex metrics from Section~2.
\end{definition}

\subsection{Analytic infrastructure}

The field-theory definition of the next section refers to several analytic
notions that Kontsevich and Segal cite as standard rather than develop: nuclear
Fréchet spaces and the trace-class maps between them, the complex manifold of
complex metrics, and holomorphic vector bundles over it. The following nodes
record the structures and properties assumed of each.

\begin{definition}[Nuclear Fréchet spaces and nuclear maps]
  \label{def:nuclear-frechet}
  \notready
  A \emph{Fréchet space} is a complete metrizable locally convex topological
  vector space over $\mathbb C$. A continuous linear map $T \colon E \to F$
  between Fréchet spaces is \emph{nuclear} (or \emph{trace-class}) if it admits
  a representation $T x = \sum_n \lambda_n\, \phi_n(x)\, y_n$ with
  $\sum_n |\lambda_n| < \infty$, the functionals $\phi_n$ equicontinuous in the
  dual $E^*$ and the vectors $y_n$ bounded in $F$. A Fréchet space is
  \emph{nuclear} if every continuous linear map from it into a Banach space is
  nuclear.

  The properties assumed of these spaces are:
  \begin{itemize}
    \item the inverse limit of a countable sequence of nuclear maps of Fréchet
      spaces is again a nuclear Fréchet space;
    \item nuclear spaces have a unique natural tensor product (the projective
      and injective tensor products coincide), coherently associative and
      commutative, so that $E \otimes F$ is unambiguous for nuclear $E$ and $F$;
    \item the strong dual of a nuclear Fréchet space is nuclear, though usually
      not metrizable, and such spaces are reflexive, so double duality recovers
      the space;
    \item a nuclear endomorphism of such a space has a well-defined trace.
  \end{itemize}

  These are standard results in the theory of nuclear spaces, cited by
  Kontsevich and Segal and not available in Mathlib. They are assumed, not
  proved.
\end{definition}

\begin{definition}[$\mathrm{Met}_{\mathbb C}(M)$ as a complex manifold]
  \label{def:metc-complex-manifold}
  \notready
  \uses{def:allowable, def:cobordism-category}
  For a smooth $d$-manifold $M$, the space $\mathrm{Met}_{\mathbb C}(M)$ of
  complex metrics on $M$, the smooth fields
  $x \mapsto g_x \in Q_{\mathbb C}(T_x M)$, is assumed to carry the structure of
  a complex manifold, in general infinite-dimensional. The complex structure on
  each fibre comes from $Q_{\mathbb C}(T_x M)$ being an open domain in the
  complex vector space of complex-valued quadratic forms on $T_x M$ (Section 2),
  and these fibrewise structures fit together holomorphically over the smoothly
  varying point $x$. Restriction of metrics along an inclusion of germs
  $U' \subset U$, and along the inclusion of a cobordism's boundary germ
  $\Sigma_i \hookrightarrow M$, are holomorphic maps
  $\mathrm{Met}_{\mathbb C}(U) \to \mathrm{Met}_{\mathbb C}(U')$ and
  $\mathrm{Met}_{\mathbb C}(M) \to \mathrm{Met}_{\mathbb C}(\Sigma_i)$.

  The notion of ``holomorphic'' used throughout Section 3 is the one attached to
  this complex manifold. Its complex-manifold structure, infinite-dimensional in
  general, is cited by Kontsevich and Segal and not available in Mathlib. It is
  assumed.
\end{definition}

\begin{definition}[Holomorphic vector bundles over $\mathrm{Met}_{\mathbb C}(M)$]
  \label{def:holomorphic-bundle}
  \notready
  \uses{def:metc-complex-manifold, def:nuclear-frechet}
  Over the complex manifold $\mathrm{Met}_{\mathbb C}(M)$ we assume the notion of
  a \emph{locally trivial holomorphic vector bundle} whose fibres are
  topological vector spaces, namely Fréchet spaces with nuclear comparison maps,
  together with:
  \begin{itemize}
    \item \emph{morphisms} of such bundles, fibrewise continuous linear maps
      depending holomorphically on the base point;
    \item \emph{pullback} of such a bundle along a holomorphic map of complex
      manifolds, in particular along the restriction maps
      $\mathrm{Met}_{\mathbb C}(M) \to \mathrm{Met}_{\mathbb C}(\Sigma_i)$ of the
      previous node.
  \end{itemize}

  This is the structure that the holomorphicity requirement on a field theory
  will later invoke: the spaces $E_{\Sigma}$ are to form such a bundle over
  $\mathrm{Met}_{\mathbb C}(U)$, and a cobordism map $Z_M$ is to be a morphism of
  such bundles over $\mathrm{Met}_{\mathbb C}(M)$, the bundles for the two ends
  being pulled back along the restriction maps. Locally trivial holomorphic
  bundles of topological vector spaces over an infinite-dimensional complex
  manifold are not available in Mathlib and are assumed.
\end{definition}
